% Created 2024-01-08 lun 09:04
% Intended LaTeX compiler: pdflatex
\documentclass[11pt]{article}
\usepackage[utf8]{inputenc}
\usepackage[T1]{fontenc}
\usepackage{graphicx}
\usepackage{longtable}
\usepackage{wrapfig}
\usepackage{rotating}
\usepackage[normalem]{ulem}
\usepackage{amsmath}
\usepackage{amssymb}
\usepackage{capt-of}
\usepackage{hyperref}
\usepackage{minted}

\usepackage{parskip}
\usepackage[utf8]{inputenc} %% For unicode chars
\usepackage{placeins}
\usepackage[margin=2.50cm]{geometry}
\usepackage[T1]{fontenc}
\usepackage{mathpazo}
\linespread{1.05}
\usepackage[scaled]{helvet}
\usepackage{courier}
\hypersetup{colorlinks=true,linkcolor=blue}
\RequirePackage{fancyvrb}
\DefineVerbatimEnvironment{verbatim}{Verbatim}{fontsize=\small,formatcom = {\color[rgb]{0.5,0,0}}}
\usepackage{mdframed}
\BeforeBeginEnvironment{minted}{\begin{mdframed}}
\AfterEndEnvironment{minted}{\end{mdframed}}
\setcounter{secnumdepth}{3}
\author{José Manuel Sánchez Álvarez}
\date{}
\title{PHP Include, Require / CRUD II}
\hypersetup{
 pdfauthor={José Manuel Sánchez Álvarez},
 pdftitle={PHP Include, Require / CRUD II},
 pdfkeywords={},
 pdfsubject={},
 pdfcreator={Emacs 28.1 (Org mode 9.5.5)}, 
 pdflang={Spanish}}
\begin{document}

\maketitle
\setcounter{tocdepth}{3}
\tableofcontents

\setlength\parindent{10pt}

\section{Update}
\label{sec:orga0c6cbd}
To add the functionality to modify (update) a car's details in your CRUD
application, you need both a PHP script to handle the update operation
and an HTML form where users can input the new car data. Below is a
step-by-step example of how you can implement this.

\subsection{Step 1: HTML Form for Car Modification}
\label{sec:orgaf7036d}
First, create an HTML form (\texttt{modify\_car.html}) where users can enter new
values for the car's attributes. For simplicity, let's assume users know
the ID of the car they want to modify.

\begin{minted}[]{html}
<!DOCTYPE html>
<html>
<head>
    <title>Modify Car Details</title>
</head>
<body>
    <h1>Modify Car Details</h1>
    <form action="update_car.php" method="post">
        Car ID (required): <input type="text" name="car_id" required><br>
        Model: <input type="text" name="model"><br>
        Price: <input type="number" name="price" step="0.01"><br>
        <!-- Add other car attributes here -->
        <input type="submit" value="Update Car">
    </form>
</body>
</html>
\end{minted}

\subsection{Step 2: PHP Script for Updating Car Details}
\label{sec:org802b95b}
Create the PHP script (\texttt{update\_car.php}) to process the form data and
update the car information in the database.

\begin{minted}[]{php}
<?php
require_once 'db_connect.php'; // Assume this file contains your PDO database connection setup

if ($_SERVER["REQUEST_METHOD"] == "POST") {
    $carId = $_POST['car_id'];
    $model = $_POST['model'] ?? null;
    $price = $_POST['price'] ?? null;
    // Retrieve other fields in a similar manner

    // Prepare the SQL statement
    $sql = "UPDATE Car SET model = :model, price = :price WHERE id = :carId";
    $stmt = $pdo->prepare($sql);

    // Bind parameters
    $stmt->bindParam(':carId', $carId, PDO::PARAM_INT);
    $stmt->bindParam(':model', $model);
    $stmt->bindParam(':price', $price);
    // Bind other parameters similarly

    try {
        $stmt->execute();
        echo "Car with ID $carId has been updated.";
    } catch (PDOException $e) {
        echo "Error updating car: " . $e->getMessage();
    }
}
?>
\end{minted}

\subsection{Explanation}
\label{sec:orgab4ce6c}
\begin{itemize}
\item \textbf{HTML Form}: The form provides inputs for the car's ID and other
attributes like model and price. The form data is sent to
\texttt{update\_car.php} upon submission.
\item \textbf{PHP Script}:

\begin{itemize}
\item The PHP script first checks if the form has been submitted.
\item It retrieves the car ID and other modified attributes from the form
data.
\item A SQL \texttt{UPDATE} statement is prepared and executed to modify the
car's details in the database.
\item The script uses PDO prepared statements to prevent SQL injection
attacks.
\end{itemize}
\end{itemize}

\subsection{Security and Error Handling}
\label{sec:orge8461f2}
\begin{itemize}
\item \textbf{Validation}: Ensure that necessary validations are done, like
checking if the car ID exists in the database.
\item \textbf{Error Handling}: Implement error handling to manage database
exceptions and form submission issues.
\item \textbf{User Feedback}: Provide clear messages to the user about the success
or failure of the update operation.
\end{itemize}

\subsection{Additional Enhancements}
\label{sec:org6db0be5}
\begin{itemize}
\item \textbf{Dynamic Data Retrieval}: For a better user experience, consider
creating a mechanism to retrieve the current details of a car (e.g.,
through a separate form or a dropdown) and populate the form with this
data for editing.
\item \textbf{Comprehensive Form Handling}: Implement more comprehensive form
validation on the server side, and consider adding client-side
validation for immediate user feedback.
\end{itemize}

By following these steps, you can effectively add the ability to modify
car details in your CRUD application, enhancing its functionality and
usability.
\end{document}
